\documentclass[b5paper,11pt]{article}
\usepackage{../notestyle}

% 文档专属：页眉文字
\fancyhead[L]{\fontsize{8pt}{10pt}\selectfont\color{gray!70}数据结构的基本概念}

% 文档信息
\title{数据结构的基本概念}
\author{ClaudeRainer}
\date{\today}

\begin{document}
\maketitle

\section{数据结构的基本概念}
\begin{cquote}
  学习编程至今已有十年，但是上次考研在08上痛失大分，不仅仅是对于概念上了解不充分，日常的工作以及开发上往往处于较高的抽象层面。
  开玩笑的说如果是在五年前的状态考，或许能够轻松拿下。

  我做过很多笔记，和数据结构相关的也不在少数。甚至发现已经完整记录过两次数据结构的笔记，笔记中记录概念的内容很多，现在来看还是没有真正理解。
  是时候重新沉淀一下，结合现在掌握的知识系统学习，以及学习如何进行“学习”。
  在如今AI快速发展的时代，如此"喧嚣"的领域，使用传统方式进行记录，再次考虑这些基础中的基础、底层下的底层，这也是极好的。

\end{cquote}

在数据结构中有四个基本概念：数据、数据元素和数据项、数据对象和数据结构、数据类型和抽象数据类型。
而组成一个数据结构则需要三个基本要素：逻辑结构、物理结构（存储结构）以及数据的运算。

\subsection{什么是数据}
数据是信息的载体，是描述客观事物属性的数值、字符以及所有能输入到计算机中并被计算机程序识别和处理的符号的集合。数据是计算机程序加工的原料。

{\color{ctp-lavender}{可以说，只要能被表示为1和0这样的基本符号的信息，均是一种数据。}}

\subsection{数据元素、数据项}
数据元素是数据的基本单位，通常作为一个整体进行考虑和处理。
一个数据元素可由若干个数据项组成，数据项（\code{原子单位}）是构成数据元素的不可分割的最小单位。

\subsection{数据结构、数据对象}
数据结构是相互之间存在一种或多种特定关系的数据元素的集合。数据结构应包括数据元素集合及元素间关系的集合，也可以说是数据的组成方式。
比如说图书馆里的书，并不只是简单的将书进行堆叠，而是按照一定的分类方法进行组织，这样就形成了一个数据结构。

\begin{adjustbox}{max width=\textwidth}
  \begin{forest} ctpTree
    [数据中心网络
      [交换机中心结构
        [树形架构]
        [平坦架构]
        [非结构性架构]
      ]
      [服务中心架构
        [巨型网络架构]
        [标准网络架构]
      ]
      [增强网络架构
        [无线网络架构]
        [光纤网络架构]
      ]
    ]
  \end{forest}
\end{adjustbox}

数据结构要研究数据元素之间的相互关系与组织方式，以及对其施加的运算及运算规则，并不关心数据元素的具体内容是什么值。
所以研究结构关系的特性而不设计其内部元素本身的内容，比如去研究数组，那么研究的是其作为数组的结构，其中的元素是什么类型或者数值是什么并不关心。

数据对象是具有相同性质的数据元素的集合，是数据的一个字集，例如整形数据对象是集合$N = \{0,1,2,3,4,5,6,7,8,9 ...\}$

\subsection{数据类型}
数据类型是一个值的集合和定义在此集合上的一组操作的总称。
\begin{itemize}
  \item 原子结构：其值不可再分的数据类型。
  \item 结构类型：其值可在再分的数据类型。
  \item 抽象数据类型：一个数学模型及定义在数学模型上的一组操作，通常是对数据的某种抽象。定义数据的取值范围及其结构形式，以及对数据操作的集合。
\end{itemize}

\section{数据结构的三要素}
数据结构有三个基本要素：逻辑结构、物理结构（存储结构）以及数据的运算。

\subsection{逻辑结构}
逻辑结构描述了数据元素之间的关系是什么，也可以说描述了数据之间的逻辑关系，从逻辑关系上描述数据。
\begin{itemize}
  \item 集合：{\color{ctp-text}{数据元素之间没有关系，数据元素之间是一个个独立的个体。}}
  \item 线性结构：{\color{ctp-text}{数据元素之间存在一对一的关系，数据元素之间存在一个线性关系。}}
  \item 树结构：{\color{ctp-text}{数据元素之间存在一对多的关系，数据元素之间存在树形结构。}}
  \item 图结构：{\color{ctp-text}{数据元素之间存在多对多的关系，数据元素之间存在图形结构。}}
\end{itemize}

% 四种逻辑结构示意图
\begin{center}
  \begin{tabular}{cc}
    % ---- 集合 ----
    \begin{tikzpicture}[baseline=(current bounding box.center)]
      \draw[dashed, color=ctp-overlay1, rounded corners=18pt]
      (-1.6,-1.0) rectangle (1.6,1.0);
      \node[circle, draw=ctp-blue, fill=ctp-blue, text=white,
      inner sep=2.5pt, font=\scriptsize] at (-0.9, 0.5) {$a_1$};
      \node[circle, draw=ctp-blue, fill=ctp-blue, text=white,
      inner sep=2.5pt, font=\scriptsize] at ( 0.3, 0.6) {$a_2$};
      \node[circle, draw=ctp-blue, fill=ctp-blue, text=white,
      inner sep=2.5pt, font=\scriptsize] at ( 1.1, 0.0) {$a_3$};
      \node[circle, draw=ctp-blue, fill=ctp-blue, text=white,
      inner sep=2.5pt, font=\scriptsize] at (-0.2,-0.5) {$a_4$};
      \node[circle, draw=ctp-blue, fill=ctp-blue, text=white,
      inner sep=2.5pt, font=\scriptsize] at (-1.1,-0.4) {$a_5$};
    \end{tikzpicture}
    &
    % ---- 线性结构 ----
    \begin{tikzpicture}[baseline=(current bounding box.center), >=latex]
      \foreach \i in {1,...,5}{
        \node[circle, draw=ctp-green, fill=ctp-green, text=white,
        inner sep=2.5pt, font=\scriptsize] (n\i) at ({(\i-1)*0.85}, 0) {$a_{\i}$};
      }
      \foreach \i [evaluate=\i as \j using int(\i+1)] in {1,...,4}{
        \draw[->, color=ctp-text, line width=0.5pt] (n\i) -- (n\j);
      }
    \end{tikzpicture}
    \\[0.4em]
    {\color{ctp-subtext0}\scriptsize 集合（无关系）}
    & {\color{ctp-subtext0}\scriptsize 线性结构（一对一）}
    \\[1.0em]
    % ---- 树结构 ----
    \begin{tikzpicture}[baseline=(current bounding box.center), >=latex]
      \node[circle, draw=ctp-mauve, fill=ctp-mauve, text=white,
      inner sep=2.5pt, font=\scriptsize] (r)  at ( 0.0, 1.0) {$a_1$};
      \node[circle, draw=ctp-teal,  fill=ctp-teal,  text=white,
      inner sep=2.5pt, font=\scriptsize] (n2) at (-1.0, 0.0) {$a_2$};
      \node[circle, draw=ctp-teal,  fill=ctp-teal,  text=white,
      inner sep=2.5pt, font=\scriptsize] (n3) at ( 1.0, 0.0) {$a_3$};
      \node[circle, draw=ctp-sky,   fill=ctp-sky,   text=white,
      inner sep=2.5pt, font=\scriptsize] (n4) at (-1.5,-1.0) {$a_4$};
      \node[circle, draw=ctp-sky,   fill=ctp-sky,   text=white,
      inner sep=2.5pt, font=\scriptsize] (n5) at (-0.5,-1.0) {$a_5$};
      \draw[->, color=ctp-text, line width=0.5pt] (r)  -- (n2);
      \draw[->, color=ctp-text, line width=0.5pt] (r)  -- (n3);
      \draw[->, color=ctp-text, line width=0.5pt] (n2) -- (n4);
      \draw[->, color=ctp-text, line width=0.5pt] (n2) -- (n5);
    \end{tikzpicture}
    &
    % ---- 图结构 ----
    \begin{tikzpicture}[baseline=(current bounding box.center)]
      \node[circle, draw=ctp-peach, fill=ctp-peach, text=white,
      inner sep=2.5pt, font=\scriptsize] (n1) at ( 0.0, 1.0) {$a_1$};
      \node[circle, draw=ctp-peach, fill=ctp-peach, text=white,
      inner sep=2.5pt, font=\scriptsize] (n2) at (-1.2, 0.2) {$a_2$};
      \node[circle, draw=ctp-peach, fill=ctp-peach, text=white,
      inner sep=2.5pt, font=\scriptsize] (n3) at ( 1.2, 0.2) {$a_3$};
      \node[circle, draw=ctp-peach, fill=ctp-peach, text=white,
      inner sep=2.5pt, font=\scriptsize] (n4) at (-0.7,-0.9) {$a_4$};
      \node[circle, draw=ctp-peach, fill=ctp-peach, text=white,
      inner sep=2.5pt, font=\scriptsize] (n5) at ( 0.7,-0.9) {$a_5$};
      \draw[color=ctp-text, line width=0.5pt] (n1) -- (n2);
      \draw[color=ctp-text, line width=0.5pt] (n1) -- (n3);
      \draw[color=ctp-text, line width=0.5pt] (n2) -- (n3);
      \draw[color=ctp-text, line width=0.5pt] (n2) -- (n4);
      \draw[color=ctp-text, line width=0.5pt] (n3) -- (n5);
      \draw[color=ctp-text, line width=0.5pt] (n4) -- (n5);
      \draw[color=ctp-text, line width=0.5pt] (n1) -- (n5);
    \end{tikzpicture}
    \\[0.4em]
    {\color{ctp-subtext0}\scriptsize 树结构（一对多）}
    & {\color{ctp-subtext0}\scriptsize 图结构（多对多）}
    \\
  \end{tabular}
\end{center}

\subsection{物理结构（存储结构）}
物理结构描述了数据元素如何存储，以及数据元素存储的顺序。
\begin{itemize}
  \item 顺序存储结构：逻辑上相邻的元素存储在物理位置上也相邻，数据元素之间的关系通过元素在存储空间中的相对位置来体现。
  \item 链式存储结构：逻辑上相邻的元素在物理位置上不一定相邻，数据元素之间的关系通过指针来体现。
  \item 散列存储结构：对数据元素的关键字施以散列函数，直接计算出该元素的存储地址，逻辑上相邻的元素物理位置不一定相邻。
  \item 索引存储结构：数据本身集中存放，另外建立附加的索引表，记录关键字到存储位置的映射，查找时先查索引表再定位数据。
\end{itemize}

以下用日常场景类比，帮助直觉上区分四种存储方式的特点：

\paragraph{顺序存储 $\leftrightarrow$ 电影院座位}
电影院的座位从第1号连续排到第$n$号，紧挨着一块儿。
要找第$k$个座位，直接用"起始位置$+$偏移量$(k-1)$"就能算出来，无需从头逐一查找——
这正是数组的随机访问原理。代价是：如果要在中间插入一排新座位，后面所有座位都得整体往后移，牵一发而动全身。

\paragraph{链式存储 $\leftrightarrow$ 寻宝游戏}
寻宝游戏里，第一张纸条告诉你"下一张在图书馆花盆下"，
到了那里又告诉你"下一张在操场大树旁"……
每个节点除存储自身数据外，还额外记录着下一个节点的内存地址（指针）。
节点可以散落在内存各处，互不相邻。
插入或删除时只需修改几个指针，不必搬运大量数据，非常灵活；
但要找第$k$个元素，必须从头顺着指针一个个走过去，效率较低。

\paragraph{散列存储 $\leftrightarrow$ 快递柜分格}
快递柜根据手机号后四位直接告诉你取件格号——不管是谁的快递，
算一下就知道放在哪格，取件时同样算一下直接去拿。
散列存储对关键字做一次哈希计算，直接得到存储地址，
查找与插入平均仅需 $O(1)$ 时间，极为高效。
代价是：哈希冲突时需要额外处理，存储空间可能有浪费，
且元素在物理上并不按逻辑顺序排列，难以进行有序遍历。

\paragraph{索引存储 $\leftrightarrow$ 书的目录}
一本书的正文按写作顺序存放，另外在开头单独建了一张目录：
"第3章$\,\rightarrow\,$第47页"。
查阅时先翻目录快速定位，再直接翻到对应页数。
索引存储也是如此：数据本身顺序摆放，另外维护一张索引表，
记录关键字到存储位置的映射。
插入新数据时需同步更新索引表，但查找时可借助索引大幅缩短搜索范围。
数据库中常见的 B$^+$ 树索引便是其典型应用。

\subsection{数据的运算}
施加在数据上的运算包括运算的定义和实现。运算的定义是针对逻辑结构的，指出运算的功能。运算的实现是针对存储结构的，指出运算的具体操作步骤。

从逻辑结构的角度举个例子，对于队列这种线性结构，需要能够实现如下的基本运算：
\begin{itemize}
  \item 入队：将一个元素加入到队列的末尾。
  \item 出队：将队列的第一个元素移除并返回。
  \item 获取队列的长度：返回队列中元素的数量。
  \item 判断队列是否为空：返回一个布尔值，表示队列是否为空。
\end{itemize}

\section{例题}
\begin{question}
  可以用（）定义一个完整的数据结构：

  \medskip
  \begin{tabular}{*{4}{p{0.22\linewidth}}}
    A.\enspace 数据元素 & B.\enspace 数据对象 & C.\enspace 数据关系 & D.\enspace 抽象数据逻辑
  \end{tabular}
\end{question}

\begin{solution}
  数据元素是数据的基本单位，一个数据元素可由若干个数据项组成，因此答案并不是\code{A}。
  而数据对象是具有相同数据元素的集合，并不能用来描述一个完整的数据结构，数据关系更是指描述数据之间联系的属于，因此答案也不是\code{B}和\code{C}。
  而抽象数据逻辑则是一个数学模型及定义在数学模型上的一组操作，通常是对数据的某种抽象。定义数据的取值范围及其结构形式，以及对数据操作的集合，
  因此可以用来定义一个完整的数据结构，所以答案是\code{D}。
\end{solution}

\begin{question}
  下列数据结构中（）是非线性数据结构

  \medskip
  \begin{tabular}{*{4}{p{0.22\linewidth}}}
    A.\enspace 树 & B.\enspace 字符串 & C.\enspace 队列 & D.\enspace 栈
  \end{tabular}
\end{question}
\begin{solution}
  树是一种非线性数据结构，因为树中的数据元素之间存在一对多的关系，数据元素之间存在树形结构。
  而字符串、队列和栈都是线性数据结构，因为它们中的数据元素之间存在一对一的关系，数据元素之间存在一个线性关系。
  因此答案是\code{A}。

  非线性数据结构，在逻辑上数据元素不是唯一前驱或唯一后继的结构。树和图是典型的非线性数据结构。
\end{solution}

\vspace{10pt} % 调整间距

\begin{question}
  下列选项中，属于逻辑结构的是（）

  \medskip
  \begin{tabular}{*{4}{p{0.22\linewidth}}}
    A.\enspace 顺序表 & B.\enspace 哈希表 & C.\enspace 有序表 & D.\enspace 单链表
  \end{tabular}
\end{question}
\begin{solution}
  顺序表、哈希表和单链表的名称本身已经蕴含了具体的存储实现方式（顺序/散列/链式），因此它们既涉及逻辑结构，也涉及物理结构。
  而有序表仅描述数据元素按值有序排列这一逻辑特性，并不限定任何存储实现方式，属于纯粹的逻辑结构描述。
  因此答案是\code{C}。

\end{solution}

\vspace{10pt} % 调整间距

\begin{question}

  下列关于数据结构的说法中，正确的是（）

  \begin{itemize}
    \item[A] 数据的逻辑结构独立于其存储结构。
    \item[B] 数据的存储结构独立于其逻辑结构。
    \item[C] 数据的逻辑结构唯一决定其存储结构。
    \item[D] 数据结构仅由其逻辑结构和存储结构决定。
  \end{itemize}
\end{question}

\begin{solution}
  数据的逻辑结构取决于面向的实际需求，采用抽象表达方式，存储结构的实际存储方式有多种不同的选择。数据的存储结构是逻辑结构在计算机上的映射，
  不能独立于逻辑结构而存在。数据的逻辑结构并不唯一决定其存储结构，数据的存储结构也不唯一决定其逻辑结构。
  数据结构还包括数据的运算，因此数据结构不仅仅由逻辑结构和存储结构决定。
  因此答案是\code{A}。
\end{solution}

\vspace{10pt} % 调整间距

\begin{question}

  在存储数据时，通常不仅要存储各数据元素的值，而且要存储（）

  \medskip
  \begin{tabular}{*{4}{p{0.22\linewidth}}}
    A.\enspace 数据的操作方法 & B.\enspace 数据元素的类型 & C.\enspace 数据元素之间的关系 & D.\enspace 数据的存取方法
  \end{tabular}
\end{question}

\begin{solution}
  在存储数据时，通常不仅要存储各数据元素的值，还要存储数据元素之间的关系。
  因为数据元素之间的关系是数据结构的一个重要组成部分，描述了数据元素之间的联系和组织方式。
  数据元素的类型虽然也很重要，但并不是所有的数据结构都需要明确存储数据元素的类型。
  数据的操作方法和存取方法则更多地涉及到算法和实现细节，而不是数据结构本身。
  因此答案是\code{C}。
\end{solution}

\vspace{10pt} % 调整间距

\begin{question}
  对于两个不同的数据结构，逻辑结构或者物理结构一定是不同的吗？
\end{question}

\begin{cquote}
  这里需要了解的是数据结构组成的三要素：逻辑结构、物理结构（存储结构）以及数据的运算。只要有一个不同就可以认为是不同的数据结构。

  逻辑结构相同，物理结构不同，最明显的例子就是线性表。物理结构可以通过数组或者链表来实现，但逻辑结构都是线性结构。

  物理数据结构相同，逻辑结构不同，可以参考栈和队列的区别。栈和队列的逻辑结构是不同的，
  栈是后进先出（LIFO）的线性结构，而队列是先进先出（FIFO）的线性结构。但是它们的物理结构可以相同，比如都可以用数组或者链表来实现。

  如果是物理结构相同，逻辑结构也相同，也有可能不是同一个数据结构，比如二叉树和二叉排序树，它们的运算定义就不相同，但物理结构和逻辑结构都是相同的。
\end{cquote}

\vspace{10pt} % 调整间距

\begin{question}
  试举一例，说明对相同的逻辑结构，同一种运算在不同的存储方式下实现时，其运算效率不同。
\end{question}
\begin{cquote}
  这很容易就想到了数组和列表两种结构，在逻辑结构中是一样的，顺序存储的方式下，线性表插入方式和删除操作的效率是$O(n)$，
  而链式存储的方式下，线性表插入方式和删除操作的效率是$O(1)$。

  而进行随机访问时，顺序存储的方式下，线性表的随机访问效率是$O(1)$，而链式存储的方式下，线性表的随机访问效率是$O(n)$。

  因此可以看出对相同的逻辑结构，同一种运算在不同的存储方式下实现时，其运算效率是不同的。
\end{cquote}

\section{总结}
从微观构成来看，数据呈现出严密的层级性。
数据（Data）是宏观的集合，是信息的载体。而深入其内部，我们会发现：数据项（Data Item）是不可分割的最小单位（如学生的学号）；
若干个数据项组合成了数据元素（Data Element），它是计算机处理数据的基本单位（如一名学生的完整记录）；
而性质相同的元素集合则构成了数据对象（Data Object）。这种“项 -> 元素 -> 对象 -> 数据”的递进关系，界定了我们处理信息的边界。

从宏观体系来看，数据结构由“三位一体”构成：逻辑结构、存储结构和数据运算。
\begin{itemize}
  \item 逻辑结构是灵魂，它脱离计算机而存在，描述的是元素之间的逻辑关系。它分为线性（一对应一）、树形（一对应多）、图形（多对应多）和集合（无关系）。
  \item 存储（物理）结构是躯壳，是逻辑结构在计算机内存中的映射。顺序存储、链式存储、索引存储和散列存储，决定了数据在内存中是“排排坐”还是“随处放”。
  \item 数据运算是桥梁，通过定义运算（在逻辑结构上）和实现运算（在存储结构上），赋予了静态数据动态处理的能力。
\end{itemize}

学习数据结构，本质上是在寻找一种平衡：如何根据问题的逻辑特性，选择合适的存储方式，从而高效地完成运算。
正如408考试常考的知识点：逻辑结构独立于存储结构，但存储结构必有其逻辑基础。
只有理解了这三者的依存关系，才能在后续学习线性表、树、图等复杂结构时游刃有余。

\end{document}
